\ifdefined\gkver\else\def\gkver{student}\fi
\documentclass[\gkver]{gaokaozhenti}
\begin{document}

\chapter{2006年江苏}

\begin{enumerate}
\item
%% number: 1
%% paperName: 2006年江苏高考第 1 题 3 分
%% typeId: 1
%% type: 单选题
%% chapter: 气体性质
%% point: 阿伏伽德罗常数
%% method: 无
%% score: 3
%% degree: 650
%% duplicateId: 0
%% body:
从下列哪一组物理量可以算出氧气的摩尔质量\xzanswer{C}

\fourchoices[answer=C]
{氧气的密度和阿伏加德罗常数}
{氧气分子的体积和阿伏加德罗常数}
{氧气分子的质量和阿伏加德罗常数}
{氧气分子的体积和氧气分子的质量}

%% answer: C
%% memo:
\memoanswer{物质的分子质量与阿伏加德罗常数的乘积即为摩尔质量，故 C 正确。}

\item
%% number: 2
%% paperName: 2006年江苏高考第 2 题 3 分
%% typeId: 1
%% type: 单选题
%% chapter: 磁场
%% point: 洛伦兹力
%% method: 无
%% score: 3
%% degree: 650
%% duplicateId: 0
%% body:
质子（p）和 $\alpha$ 粒子以相同的速率在同一匀强磁场中作匀速圆周运动，轨道半径分别为 $R_{p}$ 和 $R_{\alpha}$，周期分别为 $T_{p}$ 和 $T_{\alpha}$，则下列选项正确的是\xzanswer{A}

\fourchoices[answer=A]
{$R_{p}:R_{\alpha}=1:2$，$T_{p}:T_{\alpha}=1:2$}
{$R_{p}:R_{\alpha}=1:1$，$T_{p}:T_{\alpha}=1:1$}
{$R_{p}:R_{\alpha}=1:1$，$T_{p}:T_{\alpha}=1:2$}
{$R_{p}:R_{\alpha}=1:2$，$T_{p}:T_{\alpha}=1:1$}

%% answer: A
%% memo:
\memoanswer{由带电粒子在磁场中运动的半径公式 $R=\frac{mv}{qB}$ 得 $\frac{R_{p}}{R_{\alpha}}=\frac{m_{p}}{m_{\alpha}}\cdot\frac{q_{\alpha}}{q_{p}}=\frac{1}{4}\times\frac{2}{1}=\frac{1}{2}$；再由周期公式 $T=\frac{2\pi m}{qB}$ 得 $\frac{T_{p}}{T_{\alpha}}=\frac{m_{p}}{m_{\alpha}}\cdot\frac{q_{\alpha}}{q_{p}}=\frac{1}{4}\times\frac{2}{1}=\frac{1}{2}$，故 A 正确。}

\item
%% number: 3
%% paperName: 2006年江苏高考第 3 题 3 分
%% typeId: 1
%% type: 单选题
%% chapter: 运动和力
%% point: 牛顿定律的应用
%% method: 无
%% score: 3
%% degree: 650
%% duplicateId: 0
%% body:
一质量为 $m$ 的物体放在光滑的水平面上，今以恒力 $F$ 沿水平方向推该物体，在相同的时间间隔内，下列说法正确的是\xzanswer{D}

\fourchoices[answer=D]
{物体的位移相等}
{物体动能的变化量相等}
{$F$ 对物体做的功相等}
{物体动量的变化量相等}

%% answer: D
%% memo:
\memoanswer{由动量定理可知，物体所受冲量相等，则物体的动量变化量必然相等，故 D 正确。}

\item
%% number: 4
%% paperName: 2006年江苏高考第 4 题 3 分
%% typeId: 1
%% type: 单选题
%% chapter: 原子和原子核
%% point: 玻尔模型
%% method: 无
%% score: 3
%% degree: 650
%% duplicateId: 0
%% body:
氢原子的能级如图所示，已知可见的光的光子能量范围约为 $1.62\UeV\sim3.11\UeV$。下列说法错误的是\xzanswer{D}

\onepicture[w=5.5cm]{\includesvg{figs/04}}

\fourchoices[answer=D]
{处于 $n=3$ 能级的氢原子可以吸收任意频率的紫外线，并发生电离}
{大量氢原子从高能级向 $n=3$ 能级跃迁时，发出的光具有显著的热效应}
{大量处于 $n=4$ 能级的氢原子向低能级跃迁时，可能发出 6 种不同频率的光}
{大量处于 $n=4$ 能级的氢原子向低能级跃迁时，可能发出 3 种不同频率的可见光}

%% answer: D
%% memo:
\memoanswer{处于 $n=3$ 能级的氢原子，只需要吸收能量为 $1.51\UeV$ 的光子就能电离，故 A 正确；大量氢原子从高能级向 $n=3$ 能级跃迁时，发出的光子能量均小于 $1.51\UeV$，是红外线，具有显著的热效应，故 B 正确；大量处于 $n=4$ 能级的氢原子向低能级跃迁时，可能发出 6 种频率的光子，故 C 正确；其能量大小依次为 $12.75\UeV$（4→1）、$12.09\UeV$（3→1）、$10.2\UeV$（2→1）、$2.55\UeV$（4→2）、$1.89\UeV$（3→2）、$0.66\UeV$（4→3），其中只有两条在可见光范围，故 D 错误。}

\item
%% number: 5
%% paperName: 2006年江苏高考第 5 题 3 分
%% typeId: 1
%% type: 单选题
%% chapter: 气体性质
%% point: 热力学第一定律
%% method: 无
%% score: 3
%% degree: 650
%% duplicateId: 0
%% body:
用隔板将一绝热容器隔成 A 和 B 两部分，A 中盛有一定质量的理想气体，B 为真空（如图~\ref{2006江苏05a}）。现把隔板抽去，A 中的气体自动充满整个容器（如图~\ref{2006江苏05b}），这个过程称为气体的自由膨胀。下列说法正确的是\xzanswer{C}

\twopicture[labelA=2006江苏05a, labelB=2006江苏05b, w=6.5cm]
{figs/05a.png}
{figs/05b.png}

\fourchoices[answer=C]
{自由膨胀过程中，气体分子只作定向运动}
{自由膨胀前后，气体的压强不变}
{自由膨胀前后，气体的温度不变}
{容器中的气体在足够长的时间内，能全部自动回到 A 部分}

%% answer: C
%% memo:
\memoanswer{自由膨胀过程，气体在做定向运动的同时也在做无规则的热运动，故 A 错误；自由膨胀过程不做功，也没有热交换，故内能不变，气体的温度不变，但由于体积变大，故压强变小，B 错误、C 正确；由热力学第二定律得，自由膨胀的过程是不可逆过程，故 D 错误。}

\item
%% number: 6
%% paperName: 2006年江苏高考第 6 题 3 分
%% typeId: 1
%% type: 单选题
%% chapter: 光学
%% point: 光电效应
%% method: 无
%% score: 3
%% degree: 450
%% duplicateId: 0
%% body:
研究光电效应规律的实验装置如图所示，以频率为 $\nu$ 的光照射光电管阴极 K 时，有光电子产生。由于光电管 K、A 间加的是反向电压，光电子从阴极 K 发射后将向阳极 A 作减速运动。光电流 $I$ 由图中电流计 G 测出，反向电压 $U$ 由电压表测出。在下列表示光电效应实验规律的图象中，错误的是\xzanswer{B}

\onepicture[w=6.0cm]{figs/06.png}

\fourchoices[answer=B, ispicture=true, w=3.4cm]
{figs/06a.png}
{figs/06b.png}
{figs/06c.png}
{figs/06d.png}

%% answer: B
%% memo:
\memoanswer{由光电效应规律，截止电压 $U_{0}$ 与频率 $\nu$ 的关系满足 $eU_{0}=\frac{1}{2}mv_{m}^{2}=h\nu-W$，$U_{0}$ 应随 $\nu$ 线性增大且图线与横轴有交点（不过原点），B 图错误；由光电效应的规律可得 A、C、D 正确。}

\item
%% number: 7
%% paperName: 2006年江苏高考第 7 题 4 分
%% typeId: 2
%% type: 多选题
%% chapter: 气体性质
%% point: 热力学第一定律
%% method: 无
%% score: 4
%% degree: 650
%% duplicateId: 0
%% body:
下列说法正确的是\xzanswer{AD}

\fourchoices[answer=AD]
{气体的温度升高时，并非所有分子的速率都增大}
{盛有气体的容器作减速运动时，容器中气体的内能随之减小}
{理想气体在等容变化过程中，气体对外不做功，气体的内能不变}
{一定质量的理想气体经等温压缩后，其压强一定增大}

%% answer: AD
%% memo:
\memoanswer{气体的温度升高时，分子的平均动能增大、平均速率增大，但由于分子的热运动无规则，故并非所有分子的速率都增大，A 正确；盛有气体的容器做减速运动时，容器的动能减少，而容器中气体的温度不变，其内能不会随之减小，B 错误；理想气体在等容变化过程中，气体不对外做功，但可以进行热交换，其内能可以变化，C 错误；一定质量的理想气体经等温压缩后，其分子平均动能不变，但密度增大，单位体积内分子数增加，其压强一定增大，D 正确。}

\item
%% number: 8
%% paperName: 2006年江苏高考第 8 题 4 分
%% typeId: 2
%% type: 多选题
%% chapter: 交流电
%% point: 变压器
%% method: 无
%% score: 4
%% degree: 650
%% duplicateId: 0
%% body:
如图所示电路中的变压器为理想变压器，S 为单刀双掷开关。P 是滑动变阻器 $R$ 的滑动触头，$U_{1}$ 为加在原线圈两端的交变电压，$I_{1}$、$I_{2}$ 分别为原线圈和副线圈中的电流。下列说法正确的是\xzanswer{BC}

\onepicture[w=6.5cm]{figs/08.png}

\fourchoices[answer=BC]
{保持 P 的位置及 $U_{1}$ 不变，S 由 b 切换到 a，则 $R$ 上消耗的功率减小}
{保持 P 的位置及 $U_{1}$ 不变，S 由 a 切换到 b，则 $I_{2}$ 减小}
{保持 P 的位置及 $U_{1}$ 不变，S 由 b 切换到 a，则 $I_{1}$ 增大}
{保持 $U_{1}$ 不变，S 接在 b 端，将 P 向上滑动，则 $I_{1}$ 减小}

%% answer: BC
%% memo:
\memoanswer{由 $\frac{U_{1}}{U_{2}}=\frac{n_{1}}{n_{2}}$ 得 $P_{R}=\frac{U_{2}^{2}}{R}=\frac{n_{2}^{2}U_{1}^{2}}{n_{1}^{2}R}$，S 由 b 切换到 a 时 $n_{2}$ 增大，则 $R$ 上消耗的功率增大，故 A 错误；由 $I_{2}=\frac{U_{2}}{R}=\frac{n_{2}U_{1}}{n_{1}R}$，S 由 a 切换到 b 时 $n_{2}$ 减小，则 $I_{2}$ 减小，B 正确；同理 C 正确；由 $I_{1}=\frac{n_{2}}{n_{1}}I_{2}=\frac{n_{2}^{2}U_{1}}{n_{1}^{2}R}$，将 P 向上滑动时 $R$ 变小，则 $I_{1}$ 增大，故 D 错误。}

\item
%% number: 9
%% paperName: 2006年江苏高考第 9 题 4 分
%% typeId: 2
%% type: 多选题
%% chapter: 机械振动
%% point: 弹簧振子
%% method: 无
%% score: 4
%% degree: 650
%% duplicateId: 0
%% body:
如图所示，物体 A 置于物体 B 上，一轻质弹簧一端固定，另一端与 B 相连，在弹性限度范围内，A 和 B 一起在光滑水平面上作往复运动（不计空气阻力），并保持相对静止。则下列说法正确的是\xzanswer{AB}

\onepicture[w=4.5cm]{figs/09.png}

\fourchoices[answer=AB]
{A 和 B 均作简谐运动}
{作用在 A 上的静摩擦力大小与弹簧的形变量成正比}
{B 对 A 的静摩擦力对 A 做功，而 A 对 B 的静摩擦力对 B 不做功}
{B 对 A 的静摩擦力始终对 A 做正功，而 A 对 B 的静摩擦力始终对 B 做负功}

%% answer: AB
%% memo:
\memoanswer{将 A、B 看成整体，它们在弹簧弹力作用下一起做简谐振动，故 A 和 B 均做简谐振动，A 正确；由 $kx=(m_{A}+m_{B})a$ 及 $f_{A}=m_{A}a$ 可得 $f_{A}=\frac{m_{A}}{m_{A}+m_{B}}kx$，故 B 正确；B 对 A 的静摩擦力有时做正功、有时做负功，A 对 B 的静摩擦力也有时做正功、有时做负功，故 C、D 错误。}

\item
%% number: 10
%% paperName: 2006年江苏高考第 10 题 4 分
%% typeId: 2
%% type: 多选题
%% chapter: 功和能
%% point: 动能势能
%% method: 无
%% score: 4
%% degree: 650
%% duplicateId: 0
%% body:
江苏省沙河抽水蓄能电站自 2003 年投入运行以来，在缓解用电高峰电力紧张方面，取得了良好的社会效益和经济效益。抽水蓄能电站的工作原理是：在用电低谷时（如深夜），电站利用电网多余电能把水抽到高处蓄水池中；到用电高峰时，再利用蓄水池中的水发电。如图，蓄水池（上游水库）可视为长方体，有效总库容量（可用于发电）为 $V$，蓄水后水位高出下游水面 $H$，发电过程中上游水库水位最大落差为 $d$。统计资料表明，该电站年抽水用电为 $2.4\times10^{8}\UkWh$，年发电量为 $1.8\times10^{8}\UkWh$。则下列计算结果正确的是\xzanswer{BC}（水的密度为 $\rho$，重力加速度为 $g$，涉及重力势能的计算均以下游水面为零势能面）

\onepicture[w=5.5cm]{figs/10.png}

\fourchoices[answer=BC]
{能用于发电的水最大重力势能 $E_{v}=\rho VgH$}
{能用于发电的水的最大重力势能 $E_{p}=\rho Vg\left(H-\frac{d}{2}\right)$}
{电站的总效率达 $75\%$}
{该电站平均每天所发电能可供给一个大城市居民用电（电功率以 $10^{5}\UkW$ 计）约 $10\Uh$}

%% answer: BC
%% memo:
\memoanswer{由于上游水库的重心离下游水面的高度为 $H-\frac{d}{2}$，故能用于发电的水的最大重力势能为 $E_{p}=\rho Vg\left(H-\frac{d}{2}\right)$，A 错误、B 正确；电站的总效率 $\eta=\frac{1.8\times10^{8}}{2.4\times10^{8}}\times100\%=75\%$，C 正确；电站平均每天发电 $E_{1}=\frac{1.8\times10^{8}}{365}\UkWh=4.93\times10^{5}\UkWh$，每天可供给的用电时间约为 $4.93\Uh$，与 $10\Uh$ 相差较大，D 错误。}

\item
%% number: 11
%% paperName: 2006年江苏高考第 11 题 4 分
%% typeId: 2
%% type: 多选题
%% chapter: 机械波
%% point: 波的叠加
%% method: 无
%% score: 4
%% degree: 650
%% duplicateId: 0
%% body:
两个不等幅的脉冲波在均匀介质中均以 $1.0\Ums$ 的速率沿同一直线相向传播，$t=0$ 时刻的波形如图所示，图中小方格的边长为 $0.1\Um$。则以下不同时刻，波形正确的是\xzanswer{ABD}

\onepicture[w=8.0cm]{\includesvg{figs/11}}

\fourchoices[answer=ABD, ispicture=true, w=3.6cm]
{figs/11a.png}
{figs/11b.png}
{figs/11c.png}
{figs/11d.png}

%% answer: ABD
%% memo:
\memoanswer{由波的叠加原理，通过对两列波在相向传播过程中叠加的分析作图，可求得 A、B、D 选项正确，C 选项错误。}

\item
%% number: 12
%% paperName: 2006年江苏高考第 12 题 7 分
%% typeId: 4
%% type: 实验题
%% chapter: 光学
%% point: 测定玻璃折射率
%% method: 无
%% score: 7
%% degree: 650
%% duplicateId: 0
%% body:
\begin{enumerate}
	\item
	小球作直线运动时的频闪照片如图所示。已知频闪周期 $T=0.1\Us$，小球相邻位置间距（由照片中的刻度尺量得）分别为 $OA=6.51\Ucm$、$AB=5.59\Ucm$、$BC=4.70\Ucm$、$CD=3.80\Ucm$、$DE=2.89\Ucm$、$EF=2.00\Ucm$。小球在位置 A 时速度大小 $v_{A}=$ \tkanswer[1.5cm]{}$\Ums$，小球运动的加速度 $a=$ \tkanswer[1.5cm]{}$\Umsq$。

	\onepicture[w=10.0cm]{\includesvg{figs/12a}}

	\item
	在用插针法测定玻璃砖折射率的实验中，甲、乙、丙三位同学在纸上画出的界面 $aa'$、$bb'$ 与玻璃砖位置的关系分别如图~\ref{2006江苏12b}、~\ref{2006江苏12c} 和~\ref{2006江苏12d} 所示，其中甲、丙同学用的是矩形玻璃砖，乙同学用的是梯形玻璃砖。他们的其他操作均正确，且均以 $aa'$、$bb'$ 为界面画光路图。则甲同学测得的折射率与真实值相比 \tkanswer[1.5cm]{}（填“偏大”、“偏小”或“不变”），乙同学测得的折射率与真实值相比 \tkanswer[1.5cm]{}（填“偏大”、“偏小”或“不变”），丙同学测得的折射率与真实值相比 \tkanswer[2.5cm]{}。

	\threepicture[labelA=2006江苏12b, labelB=2006江苏12c, labelC=2006江苏12d, w=5.0cm]
	{figs/12b.png}
	{figs/12c.png}
	{figs/12d.png}
\end{enumerate}

%% answer:
\jdanswer{
\begin{enumerate}
	\item
	$0.605\Ums$；$0.9\Umsq$
	\item
	偏小；不变；可能偏大、可能偏小、可能不变
\end{enumerate}
}

%% memo:
\memoanswer{
（1）小球在位置 A 时速度大小
\[
v_{A}=\frac{OA+AB}{2T}=\frac{6.51+5.59}{2\times0.1}\times10^{-2}\Ums=0.605\Ums,
\]
由匀变速直线运动的推论 $\Delta x=aT^{2}$ 得 $CD-OA=3a_{1}T^{2}$、$DE-AB=3a_{2}T^{2}$、$EF-BC=3a_{3}T^{2}$，小球的加速度为
\[
a=\frac{a_{1}+a_{2}+a_{3}}{3}=\frac{(CD+DE+EF)-(OA+AB+BC)}{9T^{2}},
\]
代入解得 $a=-0.9\Umsq$（大小为 $0.9\Umsq$）。

（2）用图~\ref{2006江苏12b} 测定折射率时，作出的折射光线如图中虚线所示，实线表示实际光线，可见折射角增大，则由折射定律 $n=\frac{\sin i}{\sin r}$ 可知，折射率 $n$ 将偏小，故甲同学测得的折射率与真实值相比偏小；用图~\ref{2006江苏12c} 测折射率时，只要操作正确，与玻璃砖形状无关，故乙同学测得的折射率与真实值相比不变；用图~\ref{2006江苏12d} 测折射率时，折射角可能偏小、可能偏大，也可能不变，所以测得的折射率可能偏大、可能偏小，也可能不变。

\onepicture[w=6.0cm]{figs/12_an.png}
}

\item
%% number: 13
%% paperName: 2006年江苏高考第 13 题 12 分
%% typeId: 4
%% type: 实验题
%% chapter: 电路
%% point: 测电动势与内阻
%% method: 无
%% score: 12
%% degree: 450
%% duplicateId: 0
%% body:
现在按图~\ref{2006江苏13a} 所示的电路测量一节旧干电池的电动势 $E$（约 $1.5\UV$）和内阻 $r$（约 $20\UO$），可供选择的器材如下：电流表 $\mathrm{A}_{1}$、$\mathrm{A}_{2}$（量程 $0\sim500\UuA$，内阻约为 $500\UO$），滑动变阻器 $R$（阻值 $0\sim100\UO$，额定电流 $1.0\UA$），定值电阻 $R_{1}$（阻值约为 $100\UO$），电阻箱 $R_{2}$、$R_{3}$（阻值 $0\sim999.9\UO$），开关、导线若干。

由于现有电流表量程偏小，不能满足实验要求，为此，先将电流表改装（扩大量程），然后再按图~\ref{2006江苏13a} 电路进行测量。

\begin{enumerate}
	\item
	测量电流表 $\mathrm{A}_{2}$ 的内阻：按图~\ref{2006江苏13b} 电路测量 $\mathrm{A}_{2}$ 的内阻，以下给出了实验中必要的操作。
	\begin{enumerate}[label=(\Alph*)]
		\item
		断开 $\mathrm{S}_{1}$
		\item
		闭合 $\mathrm{S}_{1}$、$\mathrm{S}_{2}$
		\item
		按图~\ref{2006江苏13b} 连接线路，将滑动变阻器 $R$ 的滑片调至最左端，$R_{2}$ 调至最大
		\item
		调节 $R_{2}$，使 $\mathrm{A}_{1}$ 的示数为 $I_{1}$，记录 $R_{2}$ 的值
		\item
		断开 $\mathrm{S}_{2}$，闭合 $\mathrm{S}_{3}$
		\item
		调节滑动变阻器 $R$，使 $\mathrm{A}_{1}$、$\mathrm{A}_{2}$ 的指针偏转适中，记录 $\mathrm{A}_{1}$ 的示数 $I_{1}$
	\end{enumerate}
	请按合理顺序排列实验步骤（填序号）：\tkanswer[3cm]{}。
	\item
	将电流表 $\mathrm{A}_{2}$（较小量程）改装成电流表 A（较大量程）。如果（1）中测出 $\mathrm{A}_{2}$ 的内阻为 $468.0\UO$，现用 $R_{2}$ 将 $\mathrm{A}_{2}$ 改装成量程为 $20\UmA$ 的电流表 A，应把 $R_{2}$ 设为 \tkanswer[1.5cm]{}$\UO$ 与 $\mathrm{A}_{2}$ 并联，改装后电流表 A 的内阻 $R_{A}$ 为 \tkanswer[1.5cm]{}$\UO$。
	\item
	利用电流表 A、电阻箱 $R_{3}$ 测电池的电动势和内阻。用电流表 A、电阻箱 $R_{3}$ 及开关 S 按图~\ref{2006江苏13a} 所示电路测电池的电动势和内阻。实验时，改变 $R_{3}$ 的值，记录下电流表 A 的示数 $I$，得到若干组 $R_{3}$、$I$ 的数据，然后通过作出有关物理量的线性图象，求得电池电动势 $E$ 和内阻 $r$。
	\begin{enumerate}[label=\alph*.]
		\item
		请写出与你所作线性图象对应的函数关系式 \tkanswer[3cm]{}。
		\item
		请在图~\ref{2006江苏13c} 的虚线框内坐标中作出定性图象（要求标明两个坐标轴所代表的物理量，用符号表示）。
		\item
		图中 \tkanswer[1.5cm]{} 表示 $E$；图中 \tkanswer[1.5cm]{} 表示 $r$。
	\end{enumerate}
\end{enumerate}

\threepicture[labelA=2006江苏13a, labelB=2006江苏13b, labelC=2006江苏13c, w=4.6cm, align=right]
{figs/13a.png}
{figs/13b.png}
{\includesvg{figs/13c}}

%% answer:
\jdanswer{
\begin{enumerate}
	\item
	CBFEDA
	\item
	$12$，$11.7$
	\item
	解法一：a. $R_{3}+R_{A}=E\cdot\frac{1}{I}-r$；b. 作 $(R_{3}+R_{A})-\frac{1}{I}$ 图象；c. 直线的斜率表示 $E$，纵轴截距的绝对值表示 $r$。
	解法二：a. $\frac{1}{I}=\frac{1}{E}(R_{3}+R_{A})+\frac{r}{E}$；b. 作 $\frac{1}{I}-(R_{3}+R_{A})$ 图象；c. 直线斜率的倒数表示 $E$，纵轴截距除以斜率表示 $r$。
	解法三：a. $R_{3}=E\cdot\frac{1}{I}-(r+R_{A})$；b. 作 $R_{3}-\frac{1}{I}$ 图象；c. 直线的斜率表示 $E$，纵轴截距的绝对值与 $R_{A}$ 的差表示 $r$。三种解法的图象见详解。
\end{enumerate}
}

%% memo:
\memoanswer{
（1）由半偏法测电流表内阻的实验步骤可得其顺序为 CBFEDA。

（2）量程扩大倍数为 $n=\frac{20\times10^{-3}}{500\times10^{-6}}=40$ 倍，故并联电阻为 $R_{2}=\frac{1}{n-1}R_{g}=\frac{1}{40-1}\times468\UO=12\UO$，改装后电表的总内阻为 $R_{A}=\frac{1}{40}\times468\UO=11.7\UO$。

（3）由闭合电路的欧姆定律得 $E=I(R_{3}+R_{A}+r)$，变式可得 $R_{3}+R_{A}=E\cdot\frac{1}{I}-r$、$\frac{1}{I}=\frac{1}{E}(R_{3}+R_{A})+\frac{r}{E}$、$R_{3}=E\cdot\frac{1}{I}-(r+R_{A})$，由相应图象的斜率和截距即可求出电动势 $E$ 与内阻 $r$。

\onepicture[w=4.6cm]{figs/13_an1.png}
\onepicture[w=4.6cm]{figs/13_an2.png}
\onepicture[w=4.6cm]{figs/13_an3.png}
}

\item
%% number: 14
%% paperName: 2006年江苏高考第 14 题 14 分
%% typeId: 6
%% type: 计算题
%% chapter: 曲线运动
%% point: 人造地球卫星
%% method: 无
%% score: 14
%% degree: 650
%% duplicateId: 0
%% body:
如图所示，A 是地球的同步卫星。另一卫星 B 的圆形轨道位于赤道平面内，离地面高度为 $h$。已知地球半径为 $R$，地球自转角速度为 $\omega_{0}$，地球表面的重力加速度为 $g$，O 为地球中心。

\begin{enumerate}
	\item
	求卫星 B 的运行周期。
	\item
	如卫星 B 绕行方向与地球自转方向相同，某时刻 A、B 两卫星相距最近（O、B、A 在同一直线上），则至少经过多长时间，它们再一次相距最近？
\end{enumerate}

\onepicture[w=5.5cm, align=right]{figs/14.png}

%% answer:
\jdanswer{
\begin{enumerate}
	\item
	$T_{B}=2\pi\sqrt{\frac{(R+h)^{3}}{gR^{2}}}$
	\item
	$t=\frac{2\pi}{\sqrt{\frac{gR^{2}}{(R+h)^{3}}}-\omega_{0}}$
\end{enumerate}
}

%% memo:
\memoanswer{
（1）由万有引力定律和向心力公式得
\[
G\frac{Mm}{(R+h)^{2}}=m\frac{4\pi^{2}}{T_{B}^{2}}(R+h)\quad\text{①},
\]
\[
G\frac{Mm}{R^{2}}=mg\quad\text{②},
\]
联立①②得 $T_{B}=2\pi\sqrt{\frac{(R+h)^{3}}{gR^{2}}}$ ③。

（2）由题意和几何关系得 $(\omega_{B}-\omega_{0})t=2\pi$ ④，由③得 $\omega_{B}=\sqrt{\frac{gR^{2}}{(R+h)^{3}}}$ ⑤，代入④得 $t=\frac{2\pi}{\sqrt{\frac{gR^{2}}{(R+h)^{3}}}-\omega_{0}}$。
}

\item
%% number: 15
%% paperName: 2006年江苏高考第 15 题 14 分
%% typeId: 6
%% type: 计算题
%% chapter: 电路
%% point: 电功电功率
%% method: 无
%% score: 14
%% degree: 650
%% duplicateId: 0
%% body:
电热毯、电饭锅等是人们常用的电热式家用电器，它们一般具有加热和保温功能，其工作原理大致相同。图~\ref{2006江苏15a} 为某种电热式电器的简化电路图，主要元件有电阻丝 $R_{1}$、$R_{2}$ 和自动开关 S。

\begin{enumerate}
	\item
	当自动开关 S 闭合和断开时，用电器分别处于什么状态？
	\item
	用电器由照明电路供电（$U=220\UV$），设加热时用电器的电功率为 $400\UW$，保温时用电器的电功率为 $40\UW$，则 $R_{1}$ 和 $R_{2}$ 分别为多大？
	\item
	若将图~\ref{2006江苏15a} 中的自动开关 S 换成理想的晶体二极管 D，如图~\ref{2006江苏15b} 所示，其他条件不变，求该用电器工作 1 小时消耗的电能。
\end{enumerate}

\twopicture[labelA=2006江苏15a, labelB=2006江苏15b, w=6.5cm, align=right]
{figs/15a.png}
{figs/15b.png}

%% answer:
\jdanswer{
\begin{enumerate}
	\item
	S 闭合，处于加热状态；S 断开，处于保温状态
	\item
	$R_{1}=121\UO$，$R_{2}=1089\UO$
	\item
	$W=0.22\UkWh$（或 $7.92\times10^{5}\UJ$）
\end{enumerate}
}

%% memo:
\memoanswer{
（1）S 闭合，处于加热状态；S 断开，处于保温状态。

（2）由电功率公式得
\[
P_{1}=\frac{U^{2}}{R_{1}}\quad\text{①},\qquad P_{2}=\frac{U^{2}}{R_{1}+R_{2}}\quad\text{②},
\]
联立①②得 $R_{1}=121\UO$、$R_{2}=1089\UO$。

（3）由二极管的单向导电特性知
\[
W=P_{1}\frac{t}{2}+P_{2}\frac{t}{2}=0.22\UkWh\ (\text{或}\ 7.92\times10^{5}\UJ).
\]
}

\item
%% number: 16
%% paperName: 2006年江苏高考第 16 题 14 分
%% typeId: 6
%% type: 计算题
%% chapter: 电场
%% point: 电势能电势差电场力做功
%% method: 无
%% score: 14
%% degree: 650
%% duplicateId: 0
%% body:
如图所示，平行板电容器两极板间有场强为 $E$ 的匀强电场，且带正电的极板接地。一质量为 $m$、电荷量为 $+q$ 的带电粒子（不计重力）从 $x$ 轴上坐标为 $x_{0}$ 处静止释放。

\begin{enumerate}
	\item
	求该粒子在 $x_{0}$ 处电势能 $E_{px}$。
	\item
	试从牛顿第二定律出发，证明该带电粒子在极板间运动过程中，其动能与电势能之和保持不变。
\end{enumerate}

\onepicture[w=4.0cm, align=right]{\includesvg{figs/16}}

%% answer:
\jdanswer{
\begin{enumerate}
	\item
	$E_{px0}=-qEx_{0}$
	\item
	见解析
\end{enumerate}
}

%% memo:
\memoanswer{
（1）带电粒子运动时，电场力做的功为 $W_{\text{电}}=qEx_{0}$ ①，由功能关系知 $W_{\text{电}}=-(E_{px0}-0)$ ②，联立①②得 $E_{px0}=-qEx_{0}$ ③。

（2）解法一：在带电粒子的运动方向上任取一点，设坐标为 $x$，由牛顿第二定律得
\[
qE=ma\quad\text{④},
\]
由运动学公式得 $v_{x}^{2}=2a(x-x_{0})$ ⑤，联立④⑤可解得
\[
E_{kx}=\frac{1}{2}mv_{x}^{2}=qE(x-x_{0}),
\]
\[
E=E_{kx}+E_{px}=qE(x-x_{0})+(-qEx)=-qEx_{0},
\]
即动能与电势能之和保持不变。

解法二：在 $x$ 轴上任取两点 $x_{1}$、$x_{2}$，对应的速度大小分别为 $v_{1}$、$v_{2}$，由牛顿第二定律得 $F=qE=ma$，由运动学公式得 $v_{2}^{2}-v_{1}^{2}=2a(x_{2}-x_{1})$，联立得
\[
\frac{1}{2}mv_{2}^{2}-\frac{1}{2}mv_{1}^{2}=qE(x_{2}-x_{1}),
\]
变式得 $\frac{1}{2}mv_{2}^{2}+(-qEx_{2})=\frac{1}{2}mv_{1}^{2}+(-qEx_{1})$，即 $E_{k2}+E_{p2}=E_{k1}+E_{p1}$。
}

\item
%% number: 17
%% paperName: 2006年江苏高考第 17 题 15 分
%% typeId: 6
%% type: 计算题
%% chapter: 运动和力
%% point: 动量守恒定律
%% method: 无
%% score: 15
%% degree: 450
%% duplicateId: 0
%% body:
如图所示，质量均为 $m$ 的 A、B 两个弹性小球，用长为 $2l$ 的不可伸长的轻绳连接。现把 A、B 两球置于距地面高 $H$ 处（$H$ 足够大），间距为 $l$。当 A 球自由下落的同时，B 球以速度 $v_{0}$ 指向 A 球水平抛出。求：

\begin{enumerate}
	\item
	两球从开始运动到相碰，A 球下落的高度。
	\item
	A、B 两球碰撞（碰撞时无机械能损失）后，各自速度的水平分量。
	\item
	轻绳拉直过程中，B 球受到绳子拉力的冲量大小。
\end{enumerate}

\onepicture[w=4.5cm, align=right]{figs/17.png}

%% answer:
\jdanswer{
\begin{enumerate}
	\item
	$h=\frac{gl^{2}}{2v_{0}^{2}}$
	\item
	$v_{Ax}'=v_{0}$，$v_{Bx}'=0$
	\item
	$\frac{1}{2}mv_{0}$
\end{enumerate}
}

%% memo:
\memoanswer{
（1）设 A 球下落的高度为 $h$，由平抛规律得 $l=v_{0}t$ ①，$h=\frac{1}{2}gt^{2}$ ②，联立①②得 $h=\frac{gl^{2}}{2v_{0}^{2}}$ ③。

（2）由水平方向动量守恒得
\[
mv_{0}=mv_{Ax}'+mv_{Bx}'\quad\text{④},
\]
由机械能守恒得
\[
\frac{1}{2}m\left(v_{0}^{2}+v_{By}^{2}\right)+\frac{1}{2}mv_{Ay}^{2}=\frac{1}{2}m\left(v_{Ax}'^{2}+v_{Ay}'^{2}\right)+\frac{1}{2}m\left(v_{Bx}'^{2}+v_{By}'^{2}\right)\quad\text{⑤},
\]
式中 $v_{Ay}'=v_{Ay}$、$v_{By}'=v_{By}$，联立④⑤得 $v_{Ax}'=v_{0}$、$v_{Bx}'=0$。

（3）轻绳拉直后，两小球以相同的水平速度运动，由水平方向动量守恒得 $mv_{0}=2mv_{Bx}'$，由动量定理得 $I=mv_{Bx}'=m\frac{v_{0}}{2}$。
}

\item
%% number: 18
%% paperName: 2006年江苏高考第 18 题 15 分
%% typeId: 6
%% type: 计算题
%% chapter: 原子和原子核
%% point: 核聚变
%% method: 近似与估算
%% score: 15
%% degree: 650
%% duplicateId: 0
%% body:
天文学家测得银河系中氦的含量约为 $25\%$。有关研究表明，宇宙中氦生成的途径有两条：一是在宇宙诞生后 3 分钟左右生成的；二是在宇宙演化到恒星诞生后，由恒星内部的氢核聚变反应生成的。

\begin{enumerate}
	\item
	把氢核聚变简化为 4 个氢核（$\ce{^{1}_{1}H}$）聚变成氦核（$\ce{^{4}_{2}He}$），同时放出 2 个正电子（$\ce{^{0}_{1}e}$）和 2 个中微子（$\nu_{e}$），请写出该氢核聚变反应的方程，并计算一次反应释放的能量。
	\item
	研究表明，银河系的年龄约为 $t=3.8\times10^{17}\Us$，每秒钟银河系产生的能量约为 $1\times10^{37}\UJ$（即 $P=1\times10^{37}\UJs$）。现假定该能量全部来自上述氢核聚变反应，试估算银河系中氦的含量（最后结果保留一位有效数字）。
	\item
	根据你的估算结果，对银河系中氦的主要生成途径作出判断。
\end{enumerate}

（可能用到的数据：银河系质量约为 $M=3\times10^{41}\Ukg$，原子质量单位 $1\Uu=1.66\times10^{-27}\Ukg$，$1\Uu$ 相当于 $1.5\times10^{-10}\UJ$ 的能量，电子质量 $m_{e}=0.0005\Uu$，氦核质量 $m_{\alpha}=4.0026\Uu$，氢核质量 $m_{p}=1.0078\Uu$，中微子 $\nu_{e}$ 质量为零）

%% answer:
\jdanswer{
\begin{enumerate}
	\item
	$4\ce{^{1}_{1}H -> ^{4}_{2}He + 2^{0}_{1}e + 2\nu_{e}}$；$\Delta E=4.14\times10^{-12}\UJ$
	\item
	$2\%$
	\item
	由估算结果可知，$k\approx2\%$ 远小于 $25\%$ 的实际值，所以银河系中的氦主要是宇宙诞生后不久生成的。
\end{enumerate}
}

%% memo:
\memoanswer{
（1）氢核聚变反应的方程为 $4\ce{^{1}_{1}H -> ^{4}_{2}He + 2^{0}_{1}e + 2\nu_{e}}$。反应过程中损失的质量为 $\Delta m=4m_{p}-m_{\alpha}-2m_{e}$，由质能方程得 $\Delta E=\Delta mc^{2}=4.14\times10^{-12}\UJ$。

（2）由比例关系知，核聚变产生的氦的质量为 $m=\frac{Pt}{\Delta E}m_{\alpha}=6.1\times10^{39}\Ukg$，又银河系质量约为 $M=3\times10^{41}\Ukg$，故氦的含量为 $k=\frac{m}{M}=\frac{6.1\times10^{39}}{3\times10^{41}}\approx2\%$。

（3）由估算结果可知，$k\approx2\%$ 远小于 $25\%$ 的实际值，所以银河系中的氦主要是宇宙诞生后不久生成的。
}

\item
%% number: 19
%% paperName: 2006年江苏高考第 19 题 17 分
%% typeId: 6
%% type: 计算题
%% chapter: 电磁感应
%% point: 电磁感应能量分析
%% method: 无
%% score: 17
%% degree: 250
%% duplicateId: 0
%% body:
如图所示，顶角 $\theta=45^{\circ}$ 的金属导轨 MON 固定在水平面内，导轨处在方向竖直、磁感应强度为 $B$ 的匀强磁场中。一根与 ON 垂直的导体棒在水平外力作用下以恒定速度 $v_{0}$ 沿导轨 MON 向左滑动，导体棒的质量为 $m$，导轨与导体棒单位长度的电阻均为 $r$。导体棒与导轨接触点为 a 和 b，导体棒在滑动过程中始终保持与导轨良好接触。$t=0$ 时，导体棒位于顶角 O 处，求：

\begin{enumerate}
	\item
	$t$ 时刻流过导体棒的电流强度 $I$ 和电流方向。
	\item
	导体棒作匀速直线运动时水平外力 $F$ 的表达式。
	\item
	导体棒在 $0\sim t$ 时间内产生的焦耳热 $Q$。
	\item
	若在 $t_{0}$ 时刻将外力 $F$ 撤去，导体棒最终在导轨上静止时的坐标 $x$。
\end{enumerate}

\onepicture[w=6.0cm, align=right]{figs/19.png}

%% answer:
\jdanswer{
\begin{enumerate}
	\item
	$I=\frac{Bv_{0}}{(2+\sqrt{2})r}$，方向 $b\to a$
	\item
	$F=\frac{B^{2}v_{0}^{2}t}{(2+\sqrt{2})r}$
	\item
	$Q=\frac{B^{2}v_{0}^{3}t^{2}}{2(2+\sqrt{2})^{2}r}$
	\item
	$x=\sqrt{\frac{2(2+\sqrt{2})mv_{0}r}{B^{2}}+(v_{0}t_{0})^{2}}$
\end{enumerate}
}

%% memo:
\memoanswer{
（1）由题意知，0 到 $t$ 时间内导体棒的位移为 $x=v_{0}t$，由几何关系知 $t$ 时刻导体棒的长度为 $l=x$，导体棒切割产生的电动势为 $E=Blv_{0}$，回路的总电阻为 $R=(2x+\sqrt{2}x)r$，回路中的电流大小为
\[
I=\frac{E}{R}=\frac{Bv_{0}}{(2+\sqrt{2})r},
\]
由右手定则知，电流方向 $b\to a$。

（2）导体棒受到的安培力的大小为 $F_{\text{安}}=BlI=\frac{B^{2}v_{0}^{2}t}{(2+\sqrt{2})r}$，由导体棒的运动状态知，水平外力为 $F=F_{\text{安}}=\frac{B^{2}v_{0}^{2}t}{(2+\sqrt{2})r}$。

（3）$t$ 时刻导体棒的电功率大小为 $P=I^{2}R'=\frac{B^{2}v_{0}^{3}t}{(2+\sqrt{2})^{2}r}$，因为 $P\propto t$，利用平均功率求得产生的焦耳热为 $Q=\frac{P}{2}t=\frac{B^{2}v_{0}^{3}t^{2}}{2(2+\sqrt{2})^{2}r}$。

（4）撤去外力后，设任意时刻 $t$ 导体棒的坐标为 $x$、速度为 $v$，取很短时间 $\Delta t$ 或很短距离 $\Delta x$。在 $t\sim t+\Delta t$ 时间内，由动量定理得 $BIl\Delta t=m\Delta v$，累积求和得
\[
\sum\frac{B^{2}}{(2+\sqrt{2})r}(lv\Delta t)=\sum\frac{B^{2}}{(2+\sqrt{2})r}\Delta S=mv_{0},
\]
扫过的面积为 $\Delta S=\frac{(x_{0}+x)(x-x_{0})}{2}=\frac{x^{2}-x_{0}^{2}}{2}$（$x_{0}=v_{0}t_{0}$），解得
\[
x=\sqrt{\frac{2(2+\sqrt{2})mv_{0}r}{B^{2}}+(v_{0}t_{0})^{2}}.
\]

\onepicture[w=6.0cm]{figs/19_an.png}
}

\end{enumerate}

\end{document}
